\section* {Лекция 7. Деревья решений}

Мы успели уже рассмотреть несколько алгоритмов из класса линейных моделей, каждая из которых подходит к определённой задаче по"=своему. Однако у всех них есть один огромный недостаток: все они могли восстанавливать только линейные зависимости, т.е.
\[
y = w_0 + w_1 x_1 + w_2 x_2 + \dots
\]
а такое, на самом деле, далеко не всегда бывает в реальности. Если быть точнее, то со временем искомая зависимость в данных становится всё сложнее и сложнее, соответственно так просто воспользоваться линейными моделями мы уже не можем.

И тут на сцену выходит класс \textit{нелинейных алгоритмов}. Их суть должна быть понятна из названия "--- они способны эффективно предсказывать данные, зависимость в которых представлена нелинейно.

Можно подумать, что всё, что мы изучали ранее "--- лишь моральная подготовка к нелинейным моделям, однако в ходе занятий мы также разберёмся, почему линейные модели вообще существуют, и неужели иногда лучше строить всё же их. Тем не менее, начиная с этого занятия, мы будем изучать только нелинейные модели, и объяснять, в чём заключается их магия.